% === III. PROBLEM FORMULATION ================================================
\section{Identification Problem}
\label{sec:problem}

\subsection{Model structure}
\label{subsec:model}

The identification model is the dynamic single-track model of the baseline,
with mass $m$, yaw inertia $I_{z}$ and centre-of-gravity distances $l_{f}$,
$l_{r}$, $L=l_{f}+l_{r}$. Longitudinal and lateral dynamics are decoupled and
load transfer is neglected:
\begin{align}
\dot v_{y} &= \tfrac{1}{m}\big(\Fyr+\Fyf\cos\delta - m v_{x}\omega\big),\label{eq:vydot}\\
\dot\omega &= \tfrac{1}{I_{z}}\big(\Fyf l_{f}\cos\delta - \Fyr l_{r}\big),\label{eq:omdot}
\end{align}
with slip angles
\begin{equation}
\alpha_{f}=\delta-\arctan\!\Big(\tfrac{v_{y}+l_{f}\omega}{v_{x}}\Big),\quad
\alpha_{r}=-\arctan\!\Big(\tfrac{v_{y}-l_{r}\omega}{v_{x}}\Big).
\label{eq:slip}
\end{equation}
Axle forces follow the lateral Magic Formula
\begin{equation}
F_{y}=F_{z}\,D\sin\!\big(C\arctan\!\big(B\alpha-E(B\alpha-\arctan B\alpha)\big)\big),
\label{eq:mf}
\end{equation}
with per-axle coefficients
$\phip=[B_{f},C_{f},D_{f},E_{f},B_{r},C_{r},D_{r},E_{r}]$. The identification
state and input are $\vecx=[v_{y},\omega]$ and $\vecu=[v_{x},\delta]$, and the
nominal one-step map is $\hat\vecx_{k+1}=f(\vecx_{k},\vecu_{k},\phip)$. The
baseline discretises \eqref{eq:vydot}--\eqref{eq:omdot} by explicit Euler at
$T_{s}$; the integrator is selectable and the shipped configuration uses RK4,
because the full-scale vehicle's lateral mode is stiffer relative to $T_{s}$
than the scaled platform's. The inherited arm of Section~\ref{sec:results}
keeps Euler.

\subsection{Objective and derived quantities}

The objective is to recover $\phip$ from a single \SI{30}{\second} segment of
ordinary lap data, without dedicated manoeuvres, such that the resulting model
is both predictive and \emph{admissible} for the controller that consumes it.
Admissibility is expressed on derived quantities rather than on coefficients.
The linearised axle cornering stiffness is the slope of \eqref{eq:mf} at zero
slip,
\begin{equation}
C_{i}=F_{z,i}B_{i}C_{i}^{\mathrm{MF}}D_{i},\qquad i\in\{f,r\},
\label{eq:stiffness}
\end{equation}
(writing $C^{\mathrm{MF}}$ for the shape factor to avoid a clash), and the grip
ceiling implied by an identified set is
\begin{equation}
a_{y}^{\max}=\frac{D_{f}F_{z,f}+D_{r}F_{z,r}}{m}.
\label{eq:gripceiling}
\end{equation}
The linear model built from \eqref{eq:stiffness} is oversteering when
$l_{f}C_{f}>l_{r}C_{r}$ and is then open-loop unstable above the critical speed
$v_{\mathrm{crit}}=L\sqrt{C_{f}C_{r}/\big(m(l_{f}C_{f}-l_{r}C_{r})\big)}$
\cite{rajamani2012}.

\subsection{Practical identifiability under limited excitation}
\label{subsec:identifiability}

Two properties of \eqref{eq:mf} are used throughout. Its slope at zero
slip is $F_{z}BCD$, so the cornering stiffness depends on the coefficients only
through that product, and $E$ drops out. And if the data
never approaches the peak, \eqref{eq:mf} linearises to $F_{y}\approx
F_{z}(BCD)\alpha$, leaving $B$, $C$ and $D$ mutually unidentifiable. Given data
that does reach the peak, \eqref{eq:mf} is \emph{structurally}
identifiable~\cite{bellman1970structural}; what fails here is \emph{practical}
identifiability~\cite{raue2009practical}, the standard persistent-excitation
condition~\cite{ljung1999sysid,astrom1995adaptive} written on this model, and we
claim neither result as novel. Our claim, demonstrated in
Section~\ref{subsec:excitation}, is that the baseline pipeline's virtual-data
step can silently violate that condition, and that a bounded optimiser then
resolves the remaining freedom by sliding onto a box edge. The violation
surfaces as a converged fit rather than a failure.

\subsection{Baseline pipeline}
\label{subsec:baseline}

The baseline repeats four stages for $N_{\mathrm{it}}$ iterations
\cite{ontrack2025}. A small network is trained on the recorded lap to predict
the nominal model's one-step error from $[v_{x},v_{y},\omega,\delta]_{k}$
(stage~1); the corrected model is integrated open-loop at constant speed while
the steering angle is ramped from $0$ to $\delta_{\mathrm{sweep}}$, producing a
synthetic quasi-steady dataset (stage~2); assuming $\dot v_{y}=\dot\omega=0$ on
that data, axle forces are recovered from the yaw balance,
\begin{equation}
\Fyr=\frac{m l_{f}v_{x}\omega}{L},\qquad
\Fyf=\frac{m l_{r} v_{x}\omega}{L\cos\delta},
\label{eq:ssforces}
\end{equation}
and $\phip$ is fitted to $(\alpha_{i},F_{y,i})$ by bounded least squares
(stage~3); the fitted $\phip$ then becomes the next iteration's nominal model
(stage~4). Algorithm~\ref{alg:loop} gives the corrected cycle in full.

The rest of this paper follows from one property of that structure: stage~3
never sees the recorded lap, only the output of stage~2. The recorded lap
enters only through the residual network, so the excitation available to the
fit is set by the rollout configuration and not by how the vehicle was driven.
